\section{Review Exercises}

\subsection{Elementary computational exercises}

\begin{enumerate}
\item Compute the first six terms and determine the limit of each sequence:
\[
a_n=\frac1n,
\qquad
b_n=2-\frac1n,
\qquad
c_n=\left(\frac12\right)^n.
\]
State whether each sequence is increasing or decreasing and whether it is bounded.

\item For the sequence
\[
a_n=\frac1n,
\]
find an integer $N$ such that
\[
n\ge N\quad\Longrightarrow\quad |a_n|<0.01.
\]
Then give a formula for a suitable $N$ for an arbitrary $\varepsilon>0$.

\item Compute
\[
\lim_{n\to\infty}\frac{3n+1}{n+4},
\qquad
\lim_{n\to\infty}\frac{2n^2-n}{5n^2+3},
\qquad
\lim_{n\to\infty}\frac{4}{n^2+1}.
\]

\item Let
\[
a_n=1-\frac1{2^n}.
\]
Show directly that it is increasing and bounded above by $1$, and calculate its limit.

\item For $f(x)=2x+1$, $a=3$, and $L=7$, find a value of $\delta$ in terms of $\varepsilon$ such that
\[
0<|x-3|<\delta
\quad\Longrightarrow\quad
|f(x)-7|<\varepsilon.
\]

\item Let
\[
x_n=3+\frac{(-1)^n}{n}.
\]
Compute $f(x_n)$ for $f(x)=2x+1$ and determine its limit. Use the calculation to illustrate Heine's theorem at $a=3$.

\item Compute the limits by factoring and cancellation:
\[
\lim_{x\to4}\frac{x^2-2x-8}{x^2-9x+20},
\qquad
\lim_{x\to2}\frac{x^2-4}{x-2}.
\]
State which original expressions are undefined at the limiting point.

\item Compute
\[
\lim_{x\to0}\frac{\sqrt{1+x}-1}{x}
\]
using the conjugate. Then compute
\[
\lim_{x\to0}\frac{\sin x}{x}
\]
using the standard limit from the chapter.

\item For
\[
f(x)=
\begin{cases}
x+1,&x<0,\\
2-x,&x>0,
\end{cases}
\]
compute the left-hand and right-hand limits at $0$. Decide whether the two-sided limit exists and whether any value assigned to $f(0)$ could make the function continuous.

\item Determine the horizontal or vertical asymptotes suggested by the following limits:
\[
f(x)=\frac{3x^2+x}{x^2+4},
\qquad
g(x)=\frac1{x-2}.
\]
Compute the relevant one-sided and infinite limits. Then use the intermediate value theorem to show that
\[
h(x)=x^3-x-1
\]
has a zero in $(1,2)$.
\end{enumerate}

\subsection{Conceptual and reasoning exercises}

\begin{enumerate}
\item Explain the meaning of the definition
\[
a_n\to L.
\]
What is controlled by $\varepsilon$, what is controlled by $N$, and why only the tail of the sequence matters?

\item Compare the ordinary definition of sequence convergence with the Cauchy criterion. Why does the Cauchy condition avoid naming the candidate limit?

\item Explain why the implication ``convergent sequence $\Rightarrow$ Cauchy sequence'' follows from the triangle inequality. What deeper property of $\mathbb R$ is used in the reverse implication?

\item Explain why monotonicity alone and boundedness alone do not guarantee convergence, but their combination does for real sequences.

\item Interpret the $\varepsilon$--$\delta$ definition of a function limit using the vertical input strip and horizontal output band. Why is $x=a$ excluded?

\item State Heine's theorem in words. Explain why a function limit requires every admissible sequence approaching $a$ to produce the same output limit.

\item Why can one carefully chosen sequence disprove a proposed function limit, while one sequence cannot prove the limit? Give a concrete example.

\item Distinguish
\[
f(a)
\qquad\text{from}\qquad
\lim_{x\to a}f(x).
\]
Explain why changing a single value $f(a)$ does not change the nearby limit.

\item Explain how one-sided limits diagnose jump discontinuities and vertical asymptotes. Why must the left-hand and right-hand limits agree for a finite two-sided limit to exist?

\item The intermediate value theorem proves existence without giving a formula for a solution. Explain why continuity is the essential hypothesis and why the theorem does not guarantee uniqueness.
\end{enumerate}
